\phantomsection
\subsection*{M2A. Cyser}
\addcontentsline{toc}{subsection}{M2A. Cyser}

\textit{Cyser é um melomel feito com maçãs (geralmente na forma de sidra não fermentada, a partir do mosto de maçãs de mesa ou de maçãs próprias para sidra) ou outras frutas pomáceas, como pera e marmelo.}

\textbf{Impressão Geral}: Assemelha-se a um M1 Traditional Mead, porém com um caráter perceptível de fruta pomácea. Uma ampla gama de resultados é possível, dependendo das escolhas do hidromel base e dos ingredientes adicionais, mas o produto final deve ser equilibrado, agradável e prazeroso. O caráter da fruta deve estar em harmonia com o hidromel base e não parecer somente como uma sidra.

\textbf{Aroma}: Os aromas da fruta pomácea podem variar entre delicados, florais e lembrando vinho branco, passando por uma sidra doce e fresca, até um Apple Brandy envelhecido, dependendo da força alcoólica e do dulçor. Se variedades de fruta forem declaradas, o caráter varietal deve estar perceptível. Cysers costumam apresentar uma sugestão de acidez no aroma.

O caráter de mel pode variar de sutil a intenso, dependendo da força alcoólica e dulçor, e pode expressar características florais do néctar, refletindo quaisquer variedades de mel declaradas. O buquê de fermentação deve ser equilibrado. Hidroméis mais alcoólicos podem apresentar leves notas picantes oriundas do álcool, e os carbonatados tendem a expressar mais aroma.

Os aromas provenientes da fruta devem estar em harmonia com o hidromel base.

\textbf{Aparência}: Limpidez de razoável a brilhante; mais desejável quanto maior a limpidez. A presença de bolhas, espuma, colarinho ou efervescência deve corresponder ao nível de carbonatação declarado. A cor pode variar de palha pálida a âmbar dourado profundo (a maioria fica entre amarelo e dourado), dependendo da variedade de mel e da combinação de fruta ou sidra utilizada. Quanto mais vivas e brilhantes as cores, mais desejadas.

\textbf{Sabor}: Semelhante ao Aroma quanto a característica e intensidade de fruta e mel, tornando-se mais proeminente em versões mais alcoólicas e doces. O caráter da fruta deve ser perceptível e em equilíbrio com o nível de dulçor do hidromel. Algumas frutas podem trazer sabores básicos (amargo, doce, azedo, salgado), além de suas características próprias. Muitas frutas conferem acidez e taninos naturais, que precisam estar adequadamente em equilíbrio com o hidromel base.

Dulçor residual e final de boca devem corresponder ao nível de dulçor declarado. A acidez deve ser equilibrada, macia ou vibrante, mas não agressiva, ainda que Cyser costume apresentar acidez mais alta do que a maioria dos outros hidroméis. O caráter de fruta pomácea e de fermentação típicas de sidras e peradas também é aceitável em Cyser. Taninos podem fazer um hidromel suave parecer mais seco. Características ásperas, desagradáveis ou excessivamente sulfurosas ou notas de levedura/autólise são indesejáveis.

\textbf{Sensação na Boca}: A carbonatação segue o nível declarado, podendo ir de tranquilo a espumante. O corpo geralmente aumenta com o dulçor e a força alcoólica, tipicamente de médio-baixo a alto. A percepção de álcool acompanha a força alcoólica declarada, variando de ausente a perceptível e com sensação de aquecimento. Algumas variedades de maçã conferem considerável e perceptível nível de taninos.

\textbf{Ingredientes}: Os mesmos ingredientes de um M1 Traditional Mead, com a adição de pelo menos um tipo de fruta pomácea ou obtido a partir da mistura de hidromel e sidra.

\textbf{Comentários}: Exemplares com caráter de mel relativamente baixo podem ser melhor avaliados como C3C Experimental Cider. Um Cyser com especiarias deve ser inscrito como M3C Fruit and Spice Mead. Um Cyser com adição de outra fruta que não seja pomácea deve ser inscrito como M2E Other Fruit Mead. Um cyser com ingredientes adicionais deve ser inscrito como M4F Experimental Mead.

\textbf{Instruções para Inscrição}: Os participantes devem especificar os níveis de dulçor, carbonatação e força alcoólica. Variedades de mel podem ser declaradas. A fruta deve obrigatoriamente ser declarada. Se variedades de fruta forem especificadas, espera-se que o caráter varietal correspondente esteja presente.

\textbf{Exemplos Comerciais}: Brothers Drake Pollinator, Redstone Apple Nectar, Schramm's Mead Apple Reserve, White Winter Cyser, Williamson Idunn.
